\documentclass[10pt]{article}

\usepackage{latexsym}
\usepackage{amssymb}
\usepackage{code}
\usepackage{graphicx}

\usepackage{listings}



\usepackage[english]{babel}
%\usepackage{blindtext}
%\usepackage{fontspec}
%\setmainfont{Arial}

%\renewcommand{\familydefault}{\sfdefault}
%\usepackage{blindtext}
\pagenumbering{arabic}  % Arabic page numbers GM July 2000

\usepackage{authoraftertitle}
\usepackage{fancyhdr}
% Clear the header and footer
\fancyhead{}
\fancyfoot{}
%\lhead{\includegraphics[height=0.7cm]{../logo/nyp-logo-int.png} }
\rhead{\scriptsize \MyTitle}
%\lfoot{\scriptsize Specialist Diploma in Business \& Big Data Analytics
%  \\ Copyright
 % \copyright\ Nanyang Polytechnic. All Rights Reserved.}
\rfoot{\thepage}
\pagestyle{fancy}


\renewcommand{\topfraction}{0.95}
\renewcommand{\textfraction}{0.02}
\renewcommand{\floatpagefraction}{0.95}


\bibliographystyle{plain}

%%\textwidth      150mm \textheight     210mm \oddsidemargin  -2mm \evensidemargin -2mm

\renewcommand{\baselinestretch}{0.987}

\setlength{\parskip}{0.0in}

%\input{macros-ms}

\newcommand{\Nturns}{\, \vdash_{\mbox{\scriptsize lnf}} \,}

\newcommand{\tr}[1]{}
%%\newcommand{\tr}[1]{#1}
%%\newcommand{\nottr}[1]{#1}
\newcommand{\nottr}[1]{}

%\newcommand{\implies}{\supset}
\newcommand{\clabel}[1]{\mbox{(#1)}}
\newcommand{\rat}[1]{\rightarrowtail_{#1}}
\newcommand{\arr}{\rightarrow}
\newcommand{\arrow}{\rightarrow}
\newcommand{\Arr}{\Rightarrow}
\newcommand{\atsign}{@}
\newcommand{\simparrow}[0]{\Longleftrightarrow}
\newcommand{\proparrow}[0]{\Longrightarrow}
\newcommand{\comment}[1]{}
\newcommand{\ignore}[1]{}
%\newcommand{\kl}[1]{{\bf KL:#1}} %%%{\marginpar{\sc kl}{\bf #1}}
\newcommand{\kl}[1]{}
\newcommand{\ms}[1]{{\bf MS:#1}}       %%{\marginpar{\sc ms}{\bf #1}}
%%\newcommand{\jw}[1]{\marginpar{\sc jw}{\bf #1}}
\newcommand{\pjs}[1]{}
%%\newcommand{\ms}[1]{}
\newcommand{\jw}[1]{}
%%\newcommand{\kl}[1]{}


\newcommand{\pow}{\^{}}
\newcommand{\venv}{\Delta}
\newcommand{\mleq}{\mbox{\tt leq}}
\newcommand{\mmleq}{\leq}
\newcommand{\mas}{\mbox{\tt as}}
\newcommand{\mfix}{\mu}

\newcommand{\mysection}[1]{\vspace*{-2mm}\section{#1}\vspace*{-1mm}}
\newcommand{\mysubsection}[1]{\vspace*{-1mm}\subsection{#1}\vspace*{-1mm}}

\newcounter{cnt}
\newtheorem{ex}{Example}
%\newenvironment{example}{
%        \begin{ex}\rm}%
%    {\hfill$\Box$\end{ex}}

\newenvironment{nexample}{
        \begin{ex}\rm}%
        {\end{ex}}

%%\newtheorem{rem}{Remark}
%%\renewenvironment{remark}{
%%        \begin{rem}\rm}%
%%    {\hfill$\Box$\end{rem}}

%%\newenvironment{myexample}{
%%        \begin{ex}\rm}%
%%  {\hfill $\Diamond$\end{ex}}
%\newenvironment{example}{
%        \begin{example}\rm}%
%    {\end{example}}

\newenvironment{ttline}{\begin{trivlist}\item \tt}{\end{trivlist}}
\newenvironment{ttprog}{\begin{trivlist}\small\item \tt
        \begin{tabbing}}{\end{tabbing}\end{trivlist}}


\newcommand{\figcode}[1]
        {\begin{figure*}[t]#1
        \end{figure*}}


\title{Introduction to Spark}
%\author{Kenny Zhuo Ming Lu\\
%  \multicolumn{1}{p{.7\textwidth}}{\centering
%  \emph{School of Information Technology \\ Nanyang Polytechnic \\
%    180 Ang Mo Kio Avenue 8, Singapore 569830}}
%}


\begin{document}
\maketitle \makeatactive
\thispagestyle{fancy}

%\lstset{language=scala}

%\begin{abstract}
%\end{abstract}



\section{Exercises}

\begin{enumerate}
   \item Start the Spark cluster
\begin{code}
$/opt/spark-1.5.2-bin-hadoop2-hive2-r/sbin/start-all.sh
\end{code}
%$
   \item We are going to transform some data from HDFS. 
   \item Start Hadoop. In the VM, start a terminal and run the following command
\begin{code}
$ start-dfs.sh && start-yarn.sh
\end{code} %$
   \item Check out the source codes.
           \begin{enumerate}
               \item If you find the {\tt \~{}/git/scala-bigdata-starter/} folder exists in your system, 
\begin{code}
$ cd ~/git/scala-bigdata-starter
$ git pull
\end{code}
               \item otherwise, clone it from github.
\begin{code}
$ cd ~
$ mkdir git
$ cd git 
$ git clone http://github.com/luzhuomi/scala-bigdata-starter.git
$ cd scala-bigdata-starter
\end{code}
%$
             \end{enumerate}
 \item Change to the following folder
\begin{code}
$ cd ~/git/scala-bigdata-starter/spark-examples/
\end{code}
%$

\end{enumerate}

\subsection{Word count}

\begin{enumerate}
\item Check out the wordcount example in {\tt src/main/scala/spark/examples/WordCount.scala} which should be like as follows,
\begin{code}
  object SimpleApp {
    def main(args: Array[String]) = {
      val conf = new SparkConf().setAppName("Wordcount Application")
      val sc = new SparkContext(conf)
      val textFile = sc.textFile("hdfs://localhost:9000/input/")
      val counts = textFile.flatMap(line => line.split(" "))
      .map((word:String) => (word, 1))
      .reduceByKey(_ + _)
      counts.saveAsTextFile("hdfs://localhost:9000/output/")
    }
  }
\end{code}
%
which loads the text files from the {\tt input} folder in HDFS. For each line found in the text files, we split the line by spaces into words. 
That would give us RDD of lists of words. The {\tt flatMap} collapse the inner lists. For each word in the RDD, we associate it with the number {\tt 1}.
Then we shuffle and group them by words. Finally, the {\tt 1}s are aggregated by {\tt +}. 

\item Before running the code, we need to make sure the folders are created and the data are copied into HDFS.

\begin{code}
$ hdfs dfs -mkdir /input
$ hdfs dfs -rm -r /input/TheCompleteSherlockHolmes.txt
$ hdfs dfs -rm -r /output
$ hdfs dfs -put  ~/git/scala-bigdata-starter/spark-examples/data/wordcount/\
TheCompleteSherlockHolmes.txt /input/
\end{code}
%
Note that the two {\tt hdfs dfs -rm -r} commands will fail if the files and folders do not exist, which are normal.
\item To compile the code
\begin{code}
$ sbt package
\end{code}
%$
\item To execute the code
\begin{code}
$ /opt/spark-1.5.2-bin-hadoop2.6-hive-r/bin/spark-submit  --class SimpleApp \
target/scala-2.10/spark-examples_2.10-0.1.0.jar
\end{code}
%$
\item To observe the output 
\begin{code}
$ hdfs dfs -get /output .
$ cat output/part-*
\end{code}


\end{enumerate}

\subsection{Transformation} 
We consider an example of transforming data using Spark

\begin{enumerate}
\item Check out the example in {\tt src/main/scala/spark/examples/Transform.scala}
\begin{code}
  object Transform {
    val hdfs_nn = "127.0.0.1"
    def main(args: Array[String]) = {
      val conf = new SparkConf().setAppName("ETL (Transform) Example")
      val sc   = new SparkContext(conf)
      // load the file
      val input:RDD[String] = sc.textFile(s"hdfs://${hdfs_nn}:9000/data/transform/")
      // split by spaces
      val tokenizeds:RDD[Array[String]] = input.map(line => line.split(" "))
      tokenizeds.cache()
      
      // process all the ones
      val ones = tokenizeds
      .filter(tokenized => tokenized(0) == "1")
      .map(tokenized => { 
        val x = (tokenized(1).split(":"))(1)
        val y = (tokenized(2).split(":"))(1)
        List(x,y).mkString("\t")
      })
      ones.saveAsTextFile(s"hdfs://${hdfs_nn}:9000/output/ones")

      val zeros = tokenizeds
      .filter(tokenized => tokenized(0) == "0")
      .map(tokenized => { 
        val x = (tokenized(1).split(":"))(1)
        val y = (tokenized(2).split(":"))(1)
        List(x,y).mkString("\t")
      })
      zeros.saveAsTextFile(s"hdfs://${hdfs_nn}:9000/output/zeros")
    }
  }
\end{code}
%$
The above program parses and transform a data file in the format of 
\begin{code}
<label> 0:<x-value> 1:<y-value> 
...
<label> 0:<x-value> 1:<y-value> 
\end{code} 
into two output files where in the {\tt ones} we find all the rows with label equals to {\tt 1}  and in the {\tt zeros} we find all the rows with label {\tt 0}.
The output files are in the format of 
\begin{code}
<x-value>    <y-value>
...
<x-value>    <y-value>
\end{code}
For instance, given the input file as 
\begin{code}
1 0:102 1:230
0 0:123 1:56
0 0:22  1:2
1 0:74 1:102
\end{code}

The output files in {\tt ones} will be as 
\begin{code}
102    230
74     102
\end{code}
and those in {\tt zeros} will be as
\begin{code}
123    56
22     2
\end{code}
\item Before running the code, we need to make sure the folders are created and the data are copied into HDFS.
\begin{code}
$ hdfs dfs -mkdir /data/transform/
$ hdfs dfs -rm -r /output
$ hdfs dfs -rm -r /data/transform/input.txt
$ hdfs dfs -put  ~/git/scala-bigdata-starter/spark-examples/data/transform\
 /input.txt /data/transform/
\end{code}


\item To compile
\begin{code}
$ sbt package
\end{code}
%$
\item To execute the code
\begin{code}
$ /opt/spark-1.5.2-bin-hadoop2.6-hive-r/bin/spark-submit  --class Transform \
target/scala-2.10/spark-examples_2.10-0.1.0.jar
\end{code}
%$
\item To observe the output 
\begin{code}
$ hdfs dfs -get /output .
$ cat output/ones/part-*
$ cat output/zeros/part-*
\end{code}
%$
\end{enumerate}


\subsection{Extraction}
In this section, we consider an example of extract the US addresses from an input file.
The input file contains lines of text which could be addresses. If a US address is found in the line, the line and a ``Y'' is appended to the output, otherise the line and  an ``N'' is appended to the output.
To extract the US address we use a regular expression.
\begin{enumerate}
\item Consider the source code in {\tt src/main/scala/spark/examples/Extract.scala}
\begin{code}
  object Extract {
    val opat = compile("^(.*) ([A-Za-z]{2}) ([0-9]{5})(-[0-9]{4})?$")
    val hdfs_nn = "10.1.0.1"
    // val hdfs_nn = "127.0.0.1"
    def main(args: Array[String]) = {

      opat match 
      {
        case None    => println("Pattern compilation error." )
        case Some(p) => 
        {
          val conf = new SparkConf().setAppName("ETL (Extract) Example")
          val sc   = new SparkContext(conf)
          // load the file
          val input:RDD[String] = sc.textFile(s"hdfs://${hdfs_nn}:9000/data/extract/")

          val extracted = input.map(l => {
            exec(p,l.trim) match 
            {
              case Some(env) => List(l,"Y").mkString("\t")
              case None => List(l,"N").mkString("\t")
            }
          })
          extracted.saveAsTextFile(s"hdfs://${hdfs_nn}:9000/output/extracted")
        }
      } 
    }
  }
\end{code}
%$
\item Before running the code, we need to make sure the folders are created and the data are copied into HDFS.
\begin{code}
$ hdfs dfs -mkdir /data/extract/
$ hdfs dfs -rm -r /output
$ hdfs dfs -rm -r /data/extract/input.txt
$ hdfs dfs -put ~/git/scala-bigdata-starter/spark-examples/\
data/extract/input.txt /data/extract/
\end{code}

\item To compile
\begin{code}
$ sbt assembly
\end{code}
%$
Note that we use {\tt assembly} instead of {\tt package} due to the need of external libraries such as 
{\tt scala-pderiv} which is an efficient regular expression matching library based on partial derivative.
To ensure that the library is distributed to all the workers, we have to bundle all the depended jars into 
the main jar.

\item To execute 
\begin{code}
$ /opt/spark-1.5.2-bin-hadoop2.6-hive-r/bin/spark-submit --class Extract\ 
target/scala-2.10/spark-examples-assembly-0.1.0.jar
\end{code}
%$

\item To observe the output 
\begin{code}
$ hdfs dfs -get /output .
$ cat output/extracted/part-*
\end{code}
%$

\end{enumerate}

\end{document}
